%==============================================================================
% Appendix G: Errata and Revision History (v1 → v2)
%==============================================================================
\section{Errata and Revision History (v1 $\to$ v2)}
\label{app:errata}

This appendix documents the corrections made between v1 and v2 of this paper,
including the origin of the error, its impact on formulas and results,
and the qualitative changes in physical conclusions.

%------------------------------------------------------------------------------
\subsection{Origin of the Error}
\label{app:errata_origin}

A sign error was identified in the Riemann tensor computation within the
symbolic engine (\texttt{DPPUv2\_engine\_core\_v3.py}).

In the Riemann tensor definition
\begin{equation}
  R(X, Y)Z = [\nabla_X, \nabla_Y]Z - \nabla_{[X,Y]}Z,
\end{equation}
the commutator term $-\nabla_{[e_c, e_d]} e_b$ requires careful treatment
of the sign convention for structure constants.
Under the DPPU convention $[e_c, e_d] = -C^{e}{}_{cd}\, e_e$, we have
\begin{equation}
  -\nabla_{[e_c, e_d]} e_b
  = -\nabla_{(-C^{e}{}_{cd}\, e_e)} e_b
  = +C^{e}{}_{cd}\, \Gamma^{a}{}_{be}\, e_a.
\end{equation}
The v1 code implemented this term with a minus sign ($-C^{e}{}_{cd}\, \Gamma^{a}{}_{be}$),
effectively applying the overall minus sign twice.
This bug affects all topologies with nonzero structure constants
($\Sthree$ and $\Nilthree$) but not $\Tthree$ ($C^{i}{}_{jk} = 0$).

The error was detected during Stage 2D of DPPUv2 development when the
Weyl scalar $C^2$ was found to be identically zero for $\Nilthree$ and
squashed $\Sthree$---a result that is geometrically impossible for
non-conformally flat manifolds.
The fix corresponds to the code version update v3.0 $\to$ v3.1
(see App.~\ref{app:version_history}).

%------------------------------------------------------------------------------
\subsection{Verification of the Fix}
\label{app:errata_verification}

The corrected code was verified against known analytical results:

\begin{table}[H]
\centering
\caption{Verification of key quantities before and after the sign fix.}
\label{tab:errata_verification}
\begin{tabular}{@{}llll@{}}
\toprule
Quantity & Before fix (v1) & After fix (v2) & Expected \\
\midrule
$\Sthree$ Ricci scalar $R_{\LC}$ & $-72/r^2$ & $+24/r^2$ & $+24/r^2$ \\
Weyl $C^2$ ($\Nilthree$) & $0$ & $\neq 0$ & $\neq 0$ \\
Weyl $C^2$ (squashed $\Sthree$, $\varepsilon \neq 0$) & $0$ & $\neq 0$ & $\neq 0$ \\
\bottomrule
\end{tabular}
\end{table}

All sanity checks described in App.~\ref{app:reproducibility}---metric compatibility,
Riemann antisymmetry (3-stage), torsion tensor verification, and
Nieh--Yan density consistency---pass with the corrected code.

%------------------------------------------------------------------------------
\subsection{Impact on Formulas}
\label{app:errata_formulas}

The effective potential has the general form
$\Veff(r) = N \cdot r \left[ A r^2 + B r + C \right]$,
where $A = V^2$ (or $4V^2$ for $\Nilthree$), $B$ is the torsion-driven
coefficient proportional to $\theta_{\NY}$, and $C$ is the geometric coefficient
determined by the topology.

\subsubsection{Nieh--Yan density}

\begin{table}[H]
\centering
\caption{Nieh--Yan density $N$ before and after the fix.
After the fix, all three variants reduce to the same topology-independent form
$N \propto -V\eta/r$ for each topology.}
\label{tab:errata_NY}
\begin{tabular}{@{}llll@{}}
\toprule
Topology & Variant & v1 & v2 \\
\midrule
\multirow{3}{*}{$\Sthree$}
  & TT   & $\displaystyle -\frac{4V\eta}{r}$           & $\displaystyle -\frac{4V\eta}{r}$ (unchanged) \\[6pt]
  & REE  & $\displaystyle -\frac{2V(\eta + 4)}{r}$     & $\displaystyle -\frac{2V\eta}{r}$ \\[6pt]
  & FULL & $\displaystyle \frac{2V(4 - \eta)}{r}$      & $\displaystyle -\frac{2V\eta}{r}$ \\[6pt]
\midrule
\multirow{3}{*}{$\Nilthree$}
  & TT   & $\displaystyle -\frac{4V\eta}{r}$                    & $\displaystyle -\frac{4V\eta}{r}$ (unchanged) \\[6pt]
  & REE  & $\displaystyle \frac{2V(1 - 3\eta)}{3r}$             & $\displaystyle -\frac{2V\eta}{r}$ \\[6pt]
  & FULL & $\displaystyle -\frac{2V(3\eta + 1)}{3r}$            & $\displaystyle -\frac{2V\eta}{r}$ \\[6pt]
\bottomrule
\end{tabular}
\end{table}

A notable consequence is the identity $N_{\FULL} = N_{\REE}$ for all topologies
in v2. This is not a numerical coincidence but a structural result:
since $N_{\FULL} = N_{\TT} + N_{\REE}$ and the Levi--Civita Pontryagin density
$\varepsilon^{abcd} R^{\LC}_{abcd}$ vanishes identically on $M_3 \times \Sone$
product geometries, we have $N_{\FULL} - N_{\REE} = 0$.
In v1, the sign error introduced a spurious topology-dependent
``curvature coupling shift'' (e.g., $\eta + 4$ for $\Sthree$,
$3\eta + 1$ for $\Nilthree$) that broke this identity.

\subsubsection{Coefficient $B$ (torsion-driven term)}

\begin{table}[H]
\centering
\caption{Coefficient $B$ before and after the fix.}
\label{tab:errata_B}
\begin{tabular}{@{}llll@{}}
\toprule
Topology & Variant & v1 & v2 \\
\midrule
\multirow{3}{*}{$\Sthree$}
  & FULL & $6V\theta_{\NY}(\eta - 4)$   & $6V\eta\theta_{\NY}$ \\
  & TT   & $12V\eta\theta_{\NY}$        & $12V\eta\theta_{\NY}$ (unchanged) \\
  & REE  & $6V\theta_{\NY}(\eta + 4)$   & $6V\eta\theta_{\NY}$ \\
\midrule
\multirow{3}{*}{$\Nilthree$}
  & FULL & $8V\theta_{\NY}(3\eta + 1)$  & $24V\eta\theta_{\NY}$ \\
  & TT   & $48V\eta\theta_{\NY}$        & $48V\eta\theta_{\NY}$ (unchanged) \\
  & REE  & $8V\theta_{\NY}(3\eta - 1)$  & $24V\eta\theta_{\NY}$ \\
\bottomrule
\end{tabular}
\end{table}

In v2, all three variants have $B \propto \eta$ for each topology,
differing only in a numerical prefactor. In particular,
$B_{\TT} = 2\, B_{\FULL} = 2\, B_{\REE}$ holds universally.

\subsubsection{Coefficient $C$ (geometric term)}

\begin{table}[H]
\centering
\caption{Coefficient $C$ and its $C = 0$ roots (Type I/II boundary) before and after the fix.}
\label{tab:errata_C}
\begin{tabular}{@{}lllll@{}}
\toprule
Topology & v1 $C$ & v1 roots & v2 $C$ & v2 roots \\
\midrule
$\Sthree$   & $9(\eta + 4)^2 - 36$           & $\eta = -2,\; -6$          & $9\eta^2 - 36$  & $\eta = -2,\; +2$ \\
$\Tthree$   & $9\eta^2$                       & $\eta = 0$                 & $9\eta^2$       & $\eta = 0$ (unchanged) \\
$\Nilthree$ & $36(\eta - \tfrac{1}{3})^2 - 13$ & $\eta \approx -0.27,\; 0.94$ & $36\eta^2 + 3$  & None ($C > 0$ always) \\
\bottomrule
\end{tabular}
\end{table}

\subsubsection{Ricci scalar}

\begin{table}[H]
\centering
\caption{Ricci scalar $R$ before and after the fix.}
\label{tab:errata_ricci}
\begin{tabular}{@{}lll@{}}
\toprule
Topology & v1 & v2 \\
\midrule
$\Sthree$   & $\displaystyle \frac{2(-V^2 r^2 - 9\eta^2 - 72\eta - 108)}{3r^2}$
            & $\displaystyle \frac{2(-V^2 r^2 - 9\eta^2 + 36)}{3r^2}$ \\[8pt]
$\Tthree$   & \multicolumn{2}{c}{(unchanged)} \\[4pt]
$\Nilthree$ & $\displaystyle \frac{-4V^2 r^2 - 36\eta^2 + 24\eta + 9}{6r^2}$
            & $\displaystyle \frac{-4V^2 r^2 - 36\eta^2 - 3}{6r^2}$ \\[8pt]
\bottomrule
\end{tabular}
\end{table}

%------------------------------------------------------------------------------
\subsection{Impact on Physical Conclusions}
\label{app:errata_physics}

The corrections lead to several qualitative changes in the physical
interpretation presented in this paper.

\subsubsection{Disappearance of curvature coupling in $N$}

In v1, the Nieh--Yan density $N_{\REE}$ contained topology-dependent shifts
(e.g., $\eta + 4$ for $\Sthree$), interpreted as ``coupling'' between torsion
and background curvature. In v2, $N_{\REE} = -2V\eta/r$ is
topology-independent. The apparent curvature coupling was an artifact of
the sign error; topology dependence enters $\Veff$ solely through the
geometric coefficient $C$ via the Einstein--Hilbert term.

\subsubsection{Symmetric phase boundaries for $\Sthree$}

The Type I/II boundary (where $C = 0$) shifts from $\eta = -2,\; -6$
(asymmetric about $\eta = -4$) to $\eta = \pm 2$ (symmetric about $\eta = 0$).
This restores a natural $\eta \to -\eta$ symmetry at $\theta_{\NY} = 0$,
which is broken only when torsion ($\theta_{\NY} > 0$) is introduced via
the $B \propto \eta$ term.

\subsubsection{$\Nilthree$ coefficient $C$ is always positive}

In v1, $C_{\Nilthree}$ had a negative region ($\eta \in (-0.27,\; 0.94)$),
allowing ``narrow main band'' Type II stability driven purely by geometry.
In v2, $C_{\Nilthree} = 36\eta^2 + 3 > 0$ for all $\eta$, meaning
$\Nilthree$ can never achieve geometric-only stability without torsion.
This is consistent with the physical expectation that negative curvature
acts as a ``penalty'' for compactification.

\subsubsection{Convergence of construction variants}

In v1, the three NY variants (FULL, TT, REE) had qualitatively different
$B$ coefficients with different functional forms, leading to distinct stability
regions. The v1 discussion argued that FULL provided the ``optimal'' (widest)
stability region due to its intermediate $B$ value.

In v2, FULL and REE are identical ($N_{\FULL} = N_{\REE}$), while TT has
exactly twice the coefficient ($B_{\TT} = 2 B_{\FULL}$). The distinction
between variants is now purely quantitative (numerical prefactor), not
qualitative (functional form). The v1 claim of FULL optimality has been
revised accordingly.

%------------------------------------------------------------------------------
\subsection{Summary of Affected Sections and Figures}
\label{app:errata_sections}

\begin{table}[H]
\centering
\caption{Summary of sections and appendices modified in v2.}
\label{tab:errata_sections}
\begin{tabular}{@{}lll@{}}
\toprule
Section & Content & Nature of change \\
\midrule
Sec.~\ref{sec:reductions}  & Topology-specific reductions   & All $\Sthree$ and $\Nilthree$ formulas updated \\
Sec.~\ref{sec:numerical}   & Phase diagrams and boundaries  & Phase boundary values and tables updated \\
Sec.~\ref{sec:mechanisms}  & Type transition mechanisms      & Coefficient analysis rewritten \\
Sec.~\ref{sec:representative} & Representative points        & All $\Sthree$ and $\Nilthree$ points recalculated \\
Sec.~\ref{sec:interpretation} & Geometric interpretation     & Curvature coupling discussion revised \\
Sec.~\ref{sec:conclusions} & Conclusions                     & Summary updated to reflect corrected results \\
App.~\ref{app:theory}      & Theoretical details             & Formulas and coefficient tables corrected \\
App.~\ref{app:reproducibility} & Verification                & Updated to reflect v3.1 results \\
\bottomrule
\end{tabular}
\end{table}

All figures involving $\Sthree$ and $\Nilthree$ (11 out of 18 figures)
have been regenerated using the corrected engine (v3.1).
The five $\Tthree$ figures are unchanged.
The five $\Tthree$ figures were also regenerated using the same engine (v3.1). 
Each graph has also been revised due to changes in the representative points.

%------------------------------------------------------------------------------
\subsection{Conclusion}
\label{app:errata_conclusion}

The sign error in the Riemann tensor calculation affected all formulas involving
nonzero structure constants ($\Sthree$ and $\Nilthree$), but not $\Tthree$.
While the correction changes many quantitative results, the overall framework---the
three-type classification scheme (Type I/II/III), the role of the Nieh--Yan term
in generating torsion-driven stability, and the systematic topology dependence of
the effective potential---remains valid.

The most significant qualitative change is the simplification of the Nieh--Yan
density to a topology-independent form, which reveals a cleaner picture:
torsion provides a universal stabilization mechanism ($B \propto \eta$),
while background curvature enters independently through $C$, either lowering
the barrier ($\Sthree$, positive curvature) or raising it ($\Nilthree$,
negative curvature).
